\documentclass[9pt]{article}
\usepackage[a4paper,margin=1.45cm]{geometry}
\usepackage{amsmath,amssymb}
\usepackage{microtype}
\usepackage[hidelinks]{hyperref}
\usepackage{titlesec}
\setlength{\parindent}{0pt}
\setlength{\parskip}{2pt}
\setlength{\abovedisplayskip}{4pt}
\setlength{\belowdisplayskip}{4pt}
\setlength{\jot}{2pt}
\titleformat{\section}{\normalsize\bfseries}{}{0pt}{}
\titlespacing*{\section}{0pt}{7pt}{2pt}
\newcommand{\PE}{\operatorname{PE}}
\newcommand{\HWG}{\operatorname{HWG}}
\allowdisplaybreaks

\begin{document}
\small

\begin{center}
    {\large\bfseries Finite-coupling HWGs for $SU(3)+5F$ and $SU(3)+6F$}\\[-1mm]
    {\footnotesize General baryonic-branch formula from Bourget--Cabrera--Grimminger--Hanany--Zhong and $N_c=3$ specialisations}
\end{center}

\textbf{Notation.}
The $\mu_i$ are highest-weight fugacities for $SU(N_f)$, $\beta$ is the
finite-coupling baryonic $U(1)_B$ fugacity, and $t$ is the grading variable.
The exponent of $\beta$ is denoted $B_\beta$.  In quark-normalised baryon
number, $B(Q)=1$, one has $B=3B_\beta$ for $SU(3)$.

\section*{General finite-coupling formula}
For $N_f\geq 2N_c-1$, the finite-coupling Higgs branch consists of the
baryonic branch.  Equation~(5.52) of Bourget et al. gives
\begin{equation}
\label{eq:general-finite-hwg}
\HWG^{\mathrm{finite}}_{N_c,N_f}
=
\PE\!\left[
 t^2
 +\left(\mu_{N_c}\beta+\mu_{N_f-N_c}\beta^{-1}\right)t^{N_c}
 +\sum_{j=1}^{N_c}\mu_j\mu_{N_f-j}t^{2j}
 -\mu_{N_c}\mu_{N_f-N_c}t^{2N_c}
\right].
\end{equation}

\section*{$SU(3)+5F$: $N_c=3$, $N_f=5$}
Since $5=2\cdot3-1$, equation~\eqref{eq:general-finite-hwg} applies directly.
Before cancellation,
\begin{align}
\HWG^{\mathrm{finite}}_{5F}
=\PE\!\Big[
&t^2
+\left(\mu_3\beta+\mu_2\beta^{-1}\right)t^3
+\mu_1\mu_4t^2
+\mu_2\mu_3t^4 \notag\\
&+\mu_3\mu_2t^6
-\mu_3\mu_2t^6
\Big].
\end{align}
The two degree-six terms cancel exactly, so
\begin{align}
\HWG^{\mathrm{finite}}_{5F}
&=\PE\!\left[
(1+\mu_1\mu_4)t^2
+\left(\beta\mu_3+\beta^{-1}\mu_2\right)t^3
+\mu_2\mu_3t^4
\right] \\
&=\frac{1}
{(1-t^2)(1-\mu_1\mu_4t^2)
 (1-\beta\mu_3t^3)(1-\beta^{-1}\mu_2t^3)
 (1-\mu_2\mu_3t^4)}.
\end{align}

In $SU(5)$ Dynkin-label notation, the simplified PE input is
\begin{align}
 t^2:&\quad [0,0,0,0]_{0}+[1,0,0,1]_{0},\\
 t^3:&\quad [0,0,1,0]_{+1}+[0,1,0,0]_{-1},\\
 t^4:&\quad [0,1,1,0]_{0},
\end{align}
where the subscript is $B_\beta$.

\section*{$SU(3)+6F$: $N_c=3$, $N_f=6$}
Since $6\geq2\cdot3-1$, equation~\eqref{eq:general-finite-hwg} applies directly.
Before cancellation,
\begin{align}
\HWG^{\mathrm{finite}}_{6F}
=\PE\!\Big[
&t^2
+\mu_3\left(\beta+\beta^{-1}\right)t^3
+\mu_1\mu_5t^2
+\mu_2\mu_4t^4 \notag\\
&+\mu_3^2t^6
-\mu_3^2t^6
\Big].
\end{align}
The two degree-six terms cancel exactly, so
\begin{align}
\HWG^{\mathrm{finite}}_{6F}
&=\PE\!\left[
(1+\mu_1\mu_5)t^2
+\mu_3\left(\beta+\beta^{-1}\right)t^3
+\mu_2\mu_4t^4
\right] \\
&=\frac{1}
{(1-t^2)(1-\mu_1\mu_5t^2)
 (1-\beta\mu_3t^3)(1-\beta^{-1}\mu_3t^3)
 (1-\mu_2\mu_4t^4)}.
\end{align}

In $SU(6)$ Dynkin-label notation, the simplified PE input is
\begin{align}
 t^2:&\quad [0,0,0,0,0]_{0}+[1,0,0,0,1]_{0},\\
 t^3:&\quad [0,0,1,0,0]_{+1}+[0,0,1,0,0]_{-1},\\
 t^4:&\quad [0,1,0,1,0]_{0},
\end{align}
where the subscript is $B_\beta$.

\section*{Use as finite-coupling references}
The $5F$ result is the common finite-coupling reference for
$|k|=\frac12,\frac32,\frac52$.  The $6F$ result is the common finite-coupling
reference for $k=0,|k|=1,|k|=2$.  No instanton fugacity is introduced at
finite coupling; derived physical operator records should use $I=0$.

\vfill
\begin{thebibliography}{1}\footnotesize
\bibitem{BourgetEtAl}
A.~Bourget, S.~Cabrera, J.~F.~Grimminger, A.~Hanany and Z.~Zhong,
\emph{Brane webs and magnetic quivers for SQCD},
JHEP \textbf{03} (2020) 176,
\href{https://arxiv.org/abs/1909.00667}{arXiv:1909.00667 [hep-th]},
Eq.~(5.52).
\end{thebibliography}

\end{document}
